% Appendix A: Timeline of Major Discoveries

\chapter{Timeline of Major Discoveries}
\label{app:timeline}

This appendix consolidates the timelines from all chapters into a single chronological reference.

\section{1900--1950: Foundations}
\begin{itemize}
\item 1900: Hilbert's 23 problems (Entscheidungsproblem, consistency of arithmetic)
\item 1901: Russell's Paradox
\item 1910--13: \emph{Principia Mathematica} (Whitehead \& Russell)
\item 1928: Hilbert \& Ackermann pose Entscheidungsproblem explicitly
\item 1929: Gödel's Completeness Theorem (first-order logic)
\item 1931: Gödel's Incompleteness Theorems
\item 1936: Church ($\lambda$-calculus), Turing (Turing machines, halting problem), Post (canonical systems)
\item 1937: Turing proves TM $\equiv$ $\lambda$-calculus (PhD under Church)
\item 1943: Shannon meets Turing at Bell Labs (cryptography)
\item 1945: Von Neumann's EDVAC report (stored-program architecture)
\item 1948: Shannon: ``A Mathematical Theory of Communication''
\item 1949: Shannon: ``Communication Theory of Secrecy Systems''
\item 1950: Hamming: First error-correcting codes; Turing: ``Computing Machinery and Intelligence''
\end{itemize}

\section{1950--1980: The Classical Era}
\begin{itemize}
\item 1951: Rice's PhD thesis (Rice's Theorem)
\item 1953: Rice's Theorem published
\item 1954: Golay codes
\item 1956: Gödel's letter to von Neumann (P vs NP intuition)
\item 1959--60: BCH codes, Reed-Solomon codes
\item 1960: Solomonoff: Algorithmic probability; Rabin: Degree of difficulty; Morris: Approximate counting
\item 1965: Hartmanis-Stearns (complexity classes); Edmonds (P = tractable); Cobham (intrinsic difficulty); Flajolet-Martin (PCSA)
\item 1966: Gallager: LDPC codes; Martin-Löf: Random sequences; Rabin: Randomized Byzantine
\item 1967: Viterbi algorithm
\item 1971: Cook: SAT is NP-complete
\item 1972: Karp: 21 NP-complete problems
\item 1973: Levin: Universal search problems (USSR)
\item 1974: Fagin: Descriptive complexity (NP = $\exists$SO); Merkle: Puzzles
\item 1975: Ladner: NP-intermediate; Gray: Two Generals
\item 1976: Diffie-Hellman: Public-key crypto, key exchange; Pinski-Narin: Citation weights
\item 1977: RSA (Rivest-Shamir-Adleman); Cocks (GCHQ, classified)
\item 1978: Lamport: Logical clocks; Merkle: Thesis (Merkle trees)
\item 1979: Garey \& Johnson: \emph{Computers and Intractability}
\end{itemize}

\section{1980--2000: Expansion and Deepening}
\begin{itemize}
\item 1982: Lamport-Shostak-Pease: Byzantine Generals ($n>3t$); Goldwasser-Micali-Rackoff: ZK, IP
\item 1983: Ben-Or: Randomized consensus; Dolev-Strong: Authenticated Byzantine ($n>2t$)
\item 1985: Fischer-Lynch-Paterson: FLP Impossibility; Goldreich-Micali-Wigderson: NP $\subseteq$ CZK; ElGamal: Encryption/signatures
\item 1986: Feige-Fiat-Shamir: ZK identification, Fiat-Shamir heuristic
\item 1987: Koblitz-Miller: Elliptic Curve Cryptography; Razborov-Smolensky: AC$^0$[p] lower bounds
\item 1988: Brassard-Chaum-Crépeau: Multi-prover ZK (MIP); Dwork-Lynch-Stockmeyer: Partial synchrony
\item 1990: Flajolet et al.: LogLog; Shamir: IP = PSPACE
\item 1991: Kilian: ZK arguments, PCPs; PGP released
\item 1992: Arora-Safra: PCP Theorem; Feige-Kilian: Parallel repetition
\item 1993: Berrou et al.: Turbo codes
\item 1994: Shor: Factoring in BQP
\item 1997: Razborov-Rudich: Natural Proofs barrier; Impagliazzo-Wigderson: P = BPP from circuit bounds; GCHQ declassifies
\item 1998: Kleinberg: HITS; Brin \& Page: web search ranking
\item 1999: Castro-Liskov: PBFT
\end{itemize}

\section{2000--2020: Modern Breakthroughs}
\begin{itemize}
\item 2000: Brewer: CAP conjecture; Yao: BPP $\subseteq$ P/poly
\item 2001: Dinur: PCP theorem (alt. proof); Shokrollahi: Raptor codes
\item 2002: Gilbert-Lynch: CAP proof; Khot: Unique Games Conjecture
\item 2003: Durand-Flajolet: LogLog bias correction
\item 2004: Gyöngyi et al.: TrustRank
\item 2006: Groth-Ostrovsky-Sahai: Perfect ZK for NP; Andersen-Chung-Lang: Push PPR; Dwork-McSherry-Nissim-Smith: Differential privacy
\item 2007: Flajolet-Fusy-Gandouet-Meunier: HyperLogLog; Lamport: Fast Paxos
\item 2008: Nakamoto: Bitcoin; Groth: Short pairing-based ZK (Groth10)
\item 2009: Arıkan: Polar codes; Gentry: Fully Homomorphic Encryption
\item 2010: Google Pregel; Groth16: zk-SNARKs
\item 2013: Ben-Sasson et al.: Pinocchio; Heule-Nunkesser-Hall: HLL++ (Google)
\item 2014: Ongaro-Ousterhout: Raft; Ting: Optimal distinct elements
\item 2016: Zcash launch; Buterin-Griffith: Casper FFG
\item 2017: Tendermint: BFT PoS
\item 2018: Ben-Sasson et al.: zk-STARKs; Bünz et al.: Bulletproofs; HotStuff: Linear BFT
\item 2019: PLONK (universal setup); APPNP (PPR + GNN)
\item 2020: Halo (recursive ZK); Narwhal-Bullshark (DAG BFT); US Census 2020: differential privacy deployment
\item 2021: Kothapalli-Setty: Nova (folding); AlphaFold 2
\item 2022: Ethereum Merge (Gasper); zkEVMs (Polygon, zkSync, Scroll)
\item 2023: RISC Zero, SP1: ZK-VMs
\item 2024: NIST PQC Standards (ML-KEM, ML-DSA, SLH-DSA); AlphaProof, AlphaGeometry
\item 2025+: AI co-scientists (hypothesis generation, experiment design, formal verification)
\end{itemize}

%======================================================================
% INDEX ENTRIES
%======================================================================
\index{timeline}
\index{history of computer science}
\index{major discoveries}
\index{chronology}